% ============================================================================
% Where Does Long-Term Agent Memory Actually Fail?
% A Lifecycle-Boundary Attribution Protocol with Oracle Controls
%
% Target: ACL Rolling Review -> EMNLP / NAACL / COLING main conference (CCF-B)
%
% COMPLIANCE NOTE (venue-templates skill, checked 2026-10-18):
%   Target: ARR cycle -> EMNLP/NAACL/COLING main, long paper, initial submission
%   Primary source: https://aclrollingreview.org/cfp (re-verify at submission)
%   Content limit: 8 pages (ARR long paper); Limitations section required,
%   outside the content limit; Responsible NLP checklist at camera-ready.
%   Anonymity: two-way; the public code repository must be anonymized
%   (separate anonymous branch/artifact) before submission.
%   Official template: ACL style package (acl.sty + acl_natbib.bst) from the
%   ACL Anthology author kit -- drop the official files next to this file.
%   This skeleton is NOT the official template; it follows its structure.
%
% Compile (after adding the official ACL style files):
%   pdflatex main && bibtex main && pdflatex main && pdflatex main
% ============================================================================

\documentclass[11pt]{article}

\usepackage[hyperref]{acl}   % official ACL style package (author kit)
\usepackage{times}
\usepackage{latexsym}
\usepackage{amsmath}
\usepackage{booktabs}
\usepackage{graphicx}
\usepackage{microtype}
\usepackage{xcolor}

% Review-mode annotation helper (remove for camera-ready)
\newcommand{\todo}[1]{\textcolor{red}{[TODO: #1]}}

\newcommand{\protocol}{\textsc{Lifecycle-Boundary Attribution}}
\newcommand{\fone}{F1\textsc{\_formation}}
\newcommand{\ftwo}{F2\textsc{\_storage}}
\newcommand{\fthree}{F3\textsc{\_evolution}}
\newcommand{\ffour}{F4\textsc{\_retrieval}}
\newcommand{\ffive}{F5\textsc{\_reranking}}
\newcommand{\fsix}{F6\textsc{\_reasoning}}
\newcommand{\ezero}{E0\textsc{\_evidence\_insufficient}}
\newcommand{\evaluncertain}{\textsc{eval\_semantic\_uncertain}}
\newcommand{\rcfailure}{\textsc{rc\_reasoning\_control\_failure}}

\title{Where Does Long-Term Agent Memory Actually Fail?\\
A Lifecycle-Boundary Attribution Protocol with Oracle Controls}

\author{Anonymous submission}

\begin{document}
\maketitle

% ---------------------------------------------------------------------------
\begin{abstract}
Long-term-memory agents are evaluated end-to-end: a wrong answer is scored
wrong, and nothing is learned about \emph{where} the pipeline broke. We argue
that this measurement gap hides the dominant failure mode of memory systems,
and we fix the instrument. We present \protocol{}, a frozen protocol that
attributes every wrong answer to a specific boundary of a nine-stage memory
lifecycle (conversation $\rightarrow$ formation $\rightarrow$ storage $\rightarrow$
evolution $\rightarrow$ retrieval $\rightarrow$ reranking $\rightarrow$ context
$\rightarrow$ reasoning $\rightarrow$ answer). Three design choices make the
attribution reliable: an anti-circular evidence-sufficiency gate that compares
gold evidence against the benchmark reference answer rather than the model's
own prediction; a semantic decision layer in which lexical metrics never enter
attribution; and three orthogonal control conditions---Real Memory, Oracle
Memory (gold injected as formed memories into an independent engine), and
Oracle Context (gold injected directly into the context)---that separate
retrieval-side from reasoning-side failures. Running the protocol on two
long-term-memory benchmarks (LongMemEval-s and LoCoMo) with paired statistics
and a judge-reliability study, we find that the dominant boundary is
\emph{formation}, not retrieval: real memory loses a large fraction of gold
evidence before it ever reaches the store, and once evidence reaches the
final context the answer model rarely fails. The protocol, the measurement
validity fixes it required, and the full per-question traces are released as
an open artifact.
\end{abstract}

% ---------------------------------------------------------------------------
\section{Introduction}

Agent memory systems---retrieval-augmented stores, consolidation pipelines,
reflective memories---are increasingly the substrate of long-horizon LLM
agents, and they are evaluated almost exclusively end-to-end: answer a
question over a long conversation, score the answer. That score conflates
every way the system can fail. A missing fact, a superseded belief, a
retrieval miss, a truncated context window, and a reasoning slip all produce
the same observable: a wrong answer.

Prior work on memory failures either taxonomizes errors post-hoc from traces
or reports aggregate accuracy, and the attribution instruments underneath
carry three measurement problems that we document and fix here: (i)~lexical
scorers (exact match, token F1) are used as ground truth for correctness,
conflating representation failures with reasoning failures; (ii)~evidence
matches by lineage (a fact citing a source turn) are treated as content
preservation; (iii)~``oracle'' controls are under-specified---injecting gold
evidence into the \emph{context} and injecting it into \emph{memory} are
different interventions, and equating them misattributes retrieval failures
to reasoning.

\paragraph{Contributions.}
\begin{enumerate}
\item \textbf{A frozen attribution protocol.} \protocol{} labels every wrong
answer with the earliest lifecycle boundary at which gold evidence is not
fully present, with an evidence-sufficiency gate (\ezero{}) that is
anti-circular by construction and a semantic-judge decision layer. All
thresholds are frozen before any large-scale run (\S\ref{sec:protocol}).
\item \textbf{Three orthogonal control conditions.} Real Memory, Oracle
Memory, and Oracle Context, each measured with its own stage trace. Oracle
Memory injects gold evidence as formed memories into an independent engine
that runs the full retrieval/reranking/context pipeline; Oracle Context
bypasses it. The two controls bound where a failure can live.
\item \textbf{Measurement-validity fixes with regression tests.} Provenance
matches require content preservation; semantic evidence matching requires
faithful representation rather than topical overlap; deterministic evaluator
channels are narrowed to quote/date rules; embedding-backend failures are
loud rather than silently degrading. Each fix ships with regression tests,
including two calibration-period false negatives that pinned the bugs.
\item \textbf{An empirical study on two benchmarks} with paired statistics
(exact McNemar, paired bootstrap CIs), per-question-type breakdowns, a
cross-model replication, and an inter-judge agreement study. The dominant
failure boundary is \emph{formation}; reasoning failures (F6 candidates) are
rare and mostly control failures of the oracle conditions.
\end{enumerate}

% ---------------------------------------------------------------------------
\section{Related Work}

\paragraph{Long-term memory benchmarks.} LongMemEval \todo{cite} evaluates
chat assistants over long histories with question types that stress
extraction, temporal reasoning, and knowledge updates; LoCoMo
\todo{cite} evaluates very long multi-session conversations. We use both as
evidence-boundary substrates and treat their gold evidence spans as the
lifecycle signal.

\paragraph{Memory architectures.} RAG-style stores, consolidation and
reflection pipelines \todo{cite} optimize recall or answer accuracy. Our
contribution is orthogonal: an attribution instrument that runs on a
concrete open implementation, so the protocol's claims are checkable.

\paragraph{Failure analysis and attribution.} Error taxonomies for agents
and RAG failure analyses \todo{cite} typically require human trace reading.
Our protocol automates the lifecycle attribution while keeping the human in
the loop for judge reliability (\S\ref{sec:results}).

\paragraph{Evaluation methodology.} LLM-as-judge and semantic evaluation
\todo{cite}; judge reliability and self-preference bias \todo{cite}. We
measure inter-judge agreement explicitly (Cohen's $\kappa$) and keep the
judge model separate from the answer model in the cross-model design.

% ---------------------------------------------------------------------------
\section{Protocol: Lifecycle-Boundary Attribution}
\label{sec:protocol}

\paragraph{Lifecycle.}
Conversation $\rightarrow$ Formation $\rightarrow$ Storage $\rightarrow$
Evolution $\rightarrow$ Retrieval $\rightarrow$ Reranking $\rightarrow$
Context $\rightarrow$ Reasoning $\rightarrow$ Answer. Each arrow is a
measurable boundary with its own artifact set.

\paragraph{Labels.}
\fone{} (evidence never formed), \ftwo{} (formed but not stored),
\fthree{} (stored but superseded/dropped by governance), \ffour{} (in store
but not retrieved), \ffive{} (retrieved but dropped by reranking),
\fsix{} \textsc{\_candidate} (all evidence reached context, answer still
wrong), \ezero{} (gold evidence insufficient for the reference answer),
\evaluncertain{} (judge undecided), \rcfailure{} (oracle control answered
wrong).

\paragraph{Decision order.}
(1)~Semantic verdict on the model answer; correct $\rightarrow$ no failure.
(2)~Undecided $\rightarrow$ \evaluncertain{}. (3)~E0 gate: is the gold
evidence sufficient for the \emph{benchmark reference answer} (never the
model prediction)? Insufficient $\rightarrow$ \ezero{}. (4)~Earliest
lifecycle stage whose artifact set does not fully contain the gold evidence
$\rightarrow$ F1--F5. (5)~All stages reached and answer still wrong
$\rightarrow$ \fsix{} \textsc{\_candidate}.

\paragraph{Evidence matching tiers.}
Exact text $\rightarrow$ provenance \emph{with content-preservation check}
(shared content vocabulary $\geq 50\%$ of the smaller side and
$\geq \min(2, |gold|)$ shared tokens; rejected matches are reported as
\verb|provenance_only|) $\rightarrow$ semantic (faithful-representation
containment at 0.8, either direction, same absolute floor). Raw-token
overlap is never a match: two same-topic different-event evidence items
(``\ldots the park on 9 May 2023'' vs ``\ldots the park on 7 May 2023'')
score 0.56 under the legacy rule and are conflated, which is one mechanism
behind inflated legacy recall.

\paragraph{Unmeasured stages.}
Stages a condition does not measure carry \verb|measured: false| and are
skipped, never read as perfect recall. Oracle Context, for example, bypasses
the pipeline; its unmeasured stages are excluded from stage recall.

\paragraph{Embedding-backend discipline.}
Embeddings serve retrieval only; attribution runs on content. A failing
embedding request is retried and then fails the record loudly (hash
fallback is opt-in offline mode); inputs exceeding the backend context
window are deterministically halved; every run is verified dimension-clean
afterwards (\S\ref{sec:design}).

% ---------------------------------------------------------------------------
\section{Experimental Design}
\label{sec:design}

\paragraph{System under test.}
A concrete open implementation: a deterministic rule-based analyzer
(facts, relations, events, rules, profiles), an embedding store with dense
retrieval and reranking, context assembly, and a local answer model
(\texttt{qwen2.5:7b}, temperature 0; \texttt{gemma3:12b} cross-check),
\texttt{bge-m3} embeddings. The analyzer makes no LLM calls; only answer
generation, semantic judging, and embeddings hit the model server.

\paragraph{Conditions.}
\textbf{Real Memory}: full pipeline over the conversation.
\textbf{Oracle Memory}: gold evidence injected as formed active facts into
an independent engine, then the full retrieval/reranking/context pipeline.
\textbf{Oracle Context}: gold evidence placed directly into the answer
context, bypassing the pipeline. The three conditions partition failure
responsibility: Real vs.\ Oracle Memory isolates the memory-side loss;
Oracle Memory vs.\ Oracle Context isolates context assembly; Oracle
Context's own wrong answers are \rcfailure{}.

\paragraph{Datasets and sampling.}
LongMemEval-s (n=10 full haystack, protocol replication); LongMemEval
oracle subset (n=50, two-stage stratified by conversation and question
type, largest-remainder allocation, per-type minimum of one, fixed seed);
LoCoMo (n=50, stratified by conversation). Question-level subsets are
committed with sampling recipes; conversations are joined at load time from
the local full data; raw benchmark data is not redistributed.

\paragraph{Metrics and statistics.}
Per-condition semantic accuracy (judge verdict; lexical EM/F1 reported in
parallel, never as attribution input); per-stage recall; failure-label
distribution. Condition comparisons use exact McNemar tests on discordant
pairs and paired bootstrap percentile CIs (10,000 resamples). Judge
reliability uses Cohen's $\kappa$ between the primary judge and a second
judge model re-judging the \emph{same stored answers}.

\paragraph{Hypotheses.}
\textbf{H1}: Real Memory loses a significant fraction of gold evidence
before context; F1/F4/F5 dominate over F6 candidates.
\textbf{H2}: Oracle Memory $\gg$ Oracle Context accuracy implies
retrieval/context assembly, not reasoning, is the bottleneck.
\textbf{H3}: Real-vs-oracle inversions (oracle wrong, real right) are
control failures, not evidence of reasoning strength.

% ---------------------------------------------------------------------------
\section{Results}
\label{sec:results}

\paragraph{Study 2A (protocol replication, n=10, full haystack).}
Three-condition independence holds: \todo{real 7/10 vs oracle 10/10 vs
oracle context 10/10 (pre-fix run; to be replaced by the strict-mode
re-run)}. Calibration-period false negatives (e47becba, 58bf7951) are
explained by the formation boundary.

\paragraph{Study 2B (primary, n=50).}
\todo{Fill from results/study2b\_longmemeval\_oracle\_n50/summary.json:
semantic accuracy per condition; stage recall chain; failure distribution;
exact McNemar b/c/p; paired bootstrap CI; per-question-type table.}
Preliminary (pre-fix run): real 58.3\% vs oracle memory 72.9\% vs oracle
context 70.2\% semantic accuracy; stage recall formation 0.24, storage 1.0,
evolution 0.24, retrieval 1.0, rerank 0.99, context 0.69; failure
distribution dominated by \fone{} (14), \ezero{} (6), \evaluncertain{} (1),
zero \fsix{} candidates.

\paragraph{Study 2C (cross-benchmark, LoCoMo n=50).}
\todo{Fill: same tables; per-category breakdown; compare boundary
distribution with 2B.}

\paragraph{Study 2D (cross-model, gemma3:12b).}
\todo{Fill: same tables under the second answer model; robustness of the
formation-dominance claim.}

\paragraph{Study 2E (judge reliability).}
\todo{Fill: Cohen's $\kappa$ between qwen2.5:7b and gemma3:12b judges on a
seeded 20-record subsample of 2B (same stored answers); raw agreement;
uncertainty rate.}

% ---------------------------------------------------------------------------
\section{Analysis}

\paragraph{Where memory fails.}
The stage-recall chain localizes the loss: formation recall is the lowest
measured boundary, and context recall is high once evidence is stored. This
is an architectural statement: coverage of the extractor, not retrieval
quality, is the binding constraint for the systems under test.

\paragraph{Why end-to-end accuracy misleads.}
Lexical scoring produces false negatives on concise correct answers, which
the semantic layer resolves; conversely, lineage matches without content
preservation inflate recall. Both effects move attributed failures toward
the wrong boundary---and both are regression-tested.

\paragraph{Oracle inversions.}
Real-vs-oracle inversions are labeled \rcfailure{} and reported separately;
they bound how much of an oracle deficit is reasoning-side.

% ---------------------------------------------------------------------------
\section{Limitations}
\label{sec:limitations}

Sample sizes are n=50 per benchmark with two answer models; estimates carry
the corresponding uncertainty (reported as paired bootstrap CIs). The
analyzer is rule-based, so formation-boundary numbers characterize this
implementation family, not all memory systems. The judge is an LLM; we
quantify judge variance (Study 2E) but human adjudication of a sample
remains future work. Benchmarks are English conversational QA; temporal and
multi-hop question types are represented but not exhaustively. The public
artifact must be anonymized for review per the current call.

% ---------------------------------------------------------------------------
\section{Conclusion}

We presented \protocol{}, a frozen protocol that turns a wrong answer into
a lifecycle attribution, with three oracle controls, measurement-validity
fixes backed by regression tests, and paired-statistics reporting. On two
benchmarks and two answer models, the dominant failure boundary is
formation, not retrieval or reasoning---a conclusion invisible to
end-to-end accuracy and to attribution instruments that conflate lineage
with content or lexical with semantic correctness.

% ---------------------------------------------------------------------------
\section*{Responsible NLP Checklist}
\todo{Complete at camera-ready per ARR requirements: dataset versions and
licenses (manifest in artifact), language coverage (English), compute
(hardware, model sizes), intended use and misuse, artifact anonymization
status.}

% ---------------------------------------------------------------------------
\bibliographystyle{acl_natbib}
\bibliography{refs}

% ---------------------------------------------------------------------------
\appendix

\section{Protocol Pseudocode}
\todo{Decision-order pseudocode with the frozen thresholds.}

\section{Per-Question Attribution Tables}
\todo{Full per-question traces (appendix artifact).}

\section{Judge Prompts and Normalization}
\todo{Primary and cross-model judge prompts; normalization rules R1/R2.}

\section{Stratification Recipes}
\todo{Two-stage stratification, largest-remainder allocation, seeds.}

\section{Statistical Procedures}
\todo{Exact McNemar (binomial on discordant pairs), paired bootstrap
percentile CI, Cohen's kappa definition and estimator details.}

\section{Environment}
\todo{Models and digests, embedding backend, hardware, timings, code
commits per study.}

\end{document}
